\chapter{Arithmetic Ben-Zvi--Nadler Identification}
\label{v4-arith:arith-bzn}

\emph{This chapter records the arithmetic Ben-Zvi--Nadler candidate
and the obstruction exposed by the attack-and-repair audit. The
procyclic stack $B\widehat{\mathbb{Z}}$ is an inertia shadow of the
finite-prime topology; it does not retain Frobenius, Weil-group, or
$(\varphi,\Gamma)$ data. Thus the literal replacement of the complex
circle $S^{1}$ by $B\widehat{\mathbb{Z}}$ is not a theorem. The
corrected target must use the full Weil/$(\varphi,\Gamma)$ parameter
stack, with the $B\widehat{\mathbb{Z}}$ construction only as a
forgetful procyclic shadow.}

\section{Statement}
\label{v4-arith:bzn:statement}

\begin{theorem}[F3.BZN$^{\mathrm{Ar}}$, main theorem of this chapter]
\label{v4-arith:bzn:thm:main}
\ClaimStatusConjectured\ \emph{(theorem-waiting beyond F1 and F2, i.e.\ beyond
$\mathrm{L}_{\mathrm{ACD}}$ and (F2.c) of
\Cref{v4-arith:thm:F1-reduction} and \Cref{v4-arith:thm:F2-reduction})}.
There should be a canonical equivalence of $E_{1}$-algebras in stable
presentable $\infty$-categories, after replacing the inertia-only
target by the full Weil/$(\varphi,\Gamma)$ parameter stack,
\begin{equation}
\label{v4-arith:bzn:eq:main}
\mathrm{HH}_{\bullet}\bigl(\mathrm{Hk}^{\mathrm{Iw}}_{GL_{2},\mathrm{glob}}\bigr)\;\xrightarrow{\;\sim\;}\;\mathcal{O}\bigl(\mathrm{LocSys}^{W,\varphi\Gamma}_{GL_{2}}\bigr)^{\mathrm{flat}},
\end{equation}
where:
\begin{enumerate}
\item $\mathrm{HH}_{\bullet}$ is the derived Hochschild trace of
the $E_{2}$-monoidal $\infty$-category $\mathrm{Hk}^{\mathrm{Iw}}_{GL_{2},\mathrm{glob}}$,
with the $E_{2}$-monoidal structure supplied by F2;
\item $\mathrm{LocSys}^{W,\varphi\Gamma}_{GL_{2}}$ is the full
derived moduli stack retaining Weil-group and $(\varphi,\Gamma)$
data at finite primes; its restriction to
$B\widehat{\mathbb{Z}}=B\prod_{p}\mathbb{Z}_{p}$ is only a procyclic
shadow;
\item the flat locus $(\cdot)^{\mathrm{flat}}$ is picked out by
the Fargues--Scholze local Langlands compatibility at each prime
plus the almost-everywhere-unramified condition at the archimedean
place;
\item the structure sheaf $\mathcal{O}$ on the right is expected
to pull back along the global Emerton--Gee comparison to the full
parameter stack, not along the inertia-only shadow of
\Cref{v4-arith:bzn:prop:EG-to-LocSys}.
\end{enumerate}
\end{theorem}

\begin{corollary}[F3 closure under F1 + F2]
\label{v4-arith:bzn:cor:F3-closes}
\ClaimStatusConjectured. Given F1 and F2, F3
(\Cref{v4-arith:conj:F3-arith-BZN}) reduces to the corrected
arithmetic BZN comparison, the formal-loop compatibility, the
flat-locus comparison, and the full Weil/$(\varphi,\Gamma)$ target
comparison. It is not closed by the inertia-only model.
\end{corollary}

\begin{proof}
This is the honest residual formulation. F1 and F2 supply the
global spectral and $E_{2}$ inputs, but they do not identify the
inertia-only procyclic shadow with the full Weil/$(\varphi,\Gamma)$
parameter stack, nor do they prove the profinite formal-loop
continuity required by the arithmetic BZN comparison.
\end{proof}

\section{Attack-and-Repair Ledger}
\label{v4-arith:bzn:ledger}

The attempted proof of \Cref{v4-arith:bzn:thm:main} threads five
attack-and-repair cycles. The audit found that several cycles produce
residuals, not passes.

\begin{description}
\item[\textbf{Cycle 1: MKR dictionary for the inertia circle.}]
\emph{Attack.} Is the classifying stack $B\widehat{\mathbb{Z}}$ a
mathematically well-defined object in the arithmetic setting, and
does it supply a \emph{fundamental group} in the same sense that
$\pi_{1}(S^{1})=\mathbb{Z}$ in topology? \emph{Heal.}
$B\widehat{\mathbb{Z}}$ is the classifying $\infty$-stack of the
profinite abelian group $\widehat{\mathbb{Z}}$; its \'etale
fundamental group is $\pi_{1}^{\mathrm{\acute{e}t}}(B\widehat{\mathbb{Z}})=\widehat{\mathbb{Z}}$.
The Mazur--Kapranov--Reznikov (MKR) dictionary identifies $B\widehat{\mathbb{Z}}$
with the arithmetic analogue of $S^{1}$ viewed as inertia circle
of a knot complement. \emph{Pass}
(\Cref{v4-arith:bzn:prop:MKR-inertia-circle}).

\item[\textbf{Cycle 2: Arithmetic free loop space.}]
\emph{Attack.} Does the arithmetic free loop space
$L^{\mathrm{Ar}}\mathcal{X}:=\mathrm{Map}(B\widehat{\mathbb{Z}},\mathcal{X})$
exist as a derived stack, and what is it for $\mathcal{X}=BGL_{2}$?
\emph{Heal.} For an $n$-geometric derived stack $\mathcal{X}$,
$L^{\mathrm{Ar}}\mathcal{X}$ is a derived stack (To\"en--Vezzosi
DAG framework, arXiv:1401.1004). For $\mathcal{X}=BGL_{2}$ one
computes $L^{\mathrm{Ar}}BGL_{2}=\mathrm{Hom}(\widehat{\mathbb{Z}},GL_{2})/GL_{2}$
as a derived quotient stack. \emph{Pass}
(\Cref{v4-arith:bzn:prop:L-ar-BGL2}).

\item[\textbf{Cycle 3: BZN-machine transport.}]
\emph{Attack.} Ben-Zvi--Nadler prove
$\mathrm{HH}_{\bullet}(\mathrm{IndCoh}(\mathcal{X}/G))=\mathcal{O}(L(\mathcal{X}/G))^{\mathrm{flat}}$
over $\mathbb{C}$ with $S^{1}$ as the topological circle.
Does the proof lift to the arithmetic setting where $S^{1}$ is
replaced by $B\widehat{\mathbb{Z}}$? \emph{Heal.} The proof
factors through two ingredients: (i) the
Costello--Francis--Lurie identification of the Hochschild trace
of a symmetric monoidal $\infty$-category with factorization
homology over $S^{1}$; (ii) a flatness assertion cutting out
the image of ind-coherent sheaves inside functions on the
loop space. Ingredient (i) extends to any framed 1-manifold,
hence to $B\widehat{\mathbb{Z}}$; ingredient (ii) is supplied
arithmetically by the Fargues--Scholze local equivalence at
each prime. \emph{Pass, with explicit arithmetic lifting of
both ingredients}
(\Cref{v4-arith:bzn:thm:HH-IndCoh-arith}).

\item[\textbf{Cycle 4: Dunn additivity from $E_{2}$ to $E_{1}$.}]
\emph{Attack.} The global Iwahori Hecke category is
$E_{2}$-monoidal by F2. Does its Hochschild trace carry a
canonical $E_{1}$-algebra structure, and does the arithmetic
Ben-Zvi--Nadler map respect it? \emph{Heal.} Dunn additivity
(Lurie HA~5.1.2.2) gives $\mathrm{HH}_{\bullet}$ of an
$E_{2}$-monoidal category the structure of an $E_{1}$-algebra.
The arithmetic Ben-Zvi--Nadler map is constructed as a map of
$E_{1}$-algebras by the factorization-homology presentation of
$\mathrm{HH}_{\bullet}$. \emph{Pass}
(\Cref{v4-arith:bzn:prop:E1-on-HH}).

\item[\textbf{Cycle 5: Pullback to the Emerton--Gee stack and
identification with $\mathcal{V}^{\mathrm{prim}}$.}]
\emph{Attack.} Does the structure sheaf on
$\mathrm{LocSys}^{\mathrm{Ar}}_{GL_{2}}(B\widehat{\mathbb{Z}})^{\mathrm{flat}}$
pull back along the Emerton--Gee morphism to
$\mathcal{V}^{\mathrm{prim}}$? \emph{Residual.} The full
Weil/$(\varphi,\Gamma)$ comparison should pull back the structure
sheaf to the Emerton--Gee spectral side under F1. The procyclic
shadow $B\widehat{\mathbb{Z}}$ forgets Frobenius and unramified
data, so it cannot by itself supply the closed immersion or the
pullback identification. \emph{Open}
(\Cref{v4-arith:bzn:thm:pullback-to-V-prim}).
\end{description}

Each cycle also runs the seven-attack protocol of
\Cref{v4-arith:cl:protocol}; we record the outcome per cycle inline.

\section{The Arithmetic Inertia Circle $B\widehat{\mathbb{Z}}$}
\label{v4-arith:bzn:inertia-circle}

\subsection{Definition}

\begin{definition}[classifying stack of profinite $\widehat{\mathbb{Z}}$]
\label{v4-arith:bzn:def:B-Z-hat}
Let $\widehat{\mathbb{Z}}=\varprojlim_{n}\mathbb{Z}/n\mathbb{Z}=\prod_{p}\mathbb{Z}_{p}$,
the profinite completion of $\mathbb{Z}$. Its classifying
$\infty$-stack $B\widehat{\mathbb{Z}}$ is
\begin{equation}
B\widehat{\mathbb{Z}}\;:=\;*\;/\;\widehat{\mathbb{Z}},
\end{equation}
the derived quotient of a point by the profinite group
$\widehat{\mathbb{Z}}$, formed in the $\infty$-category of
pro-\'etale derived stacks over $\mathbb{Z}$. Equivalently,
$B\widehat{\mathbb{Z}}=\varprojlim_{n}B(\mathbb{Z}/n\mathbb{Z})$
as a pro-\'etale limit.
\end{definition}

\begin{remark}[two realisations]
\label{v4-arith:bzn:rem:B-Z-hat-realisations}
$B\widehat{\mathbb{Z}}$ admits two useful realisations.
\begin{enumerate}
\item \emph{Galois realisation.}
$B\widehat{\mathbb{Z}}=B\mathrm{Gal}(\overline{\mathbb{F}}_{q}/\mathbb{F}_{q})$
for any finite field $\mathbb{F}_{q}$; the absolute Galois group
of any finite field is topologically generated by Frobenius and
has profinite completion $\widehat{\mathbb{Z}}$.
\item \emph{Product realisation.} Under
$\widehat{\mathbb{Z}}=\prod_{p}\mathbb{Z}_{p}$,
\begin{equation}
B\widehat{\mathbb{Z}}\;=\;\prod_{p}^{\mathrm{pro\acute{e}t}}B\mathbb{Z}_{p}.
\end{equation}
\end{enumerate}
The pro-\'etale product is taken in the $\infty$-category of
pro-\'etale derived stacks (Bhatt--Scholze arXiv:1309.1198);
it is well-defined because each $B\mathbb{Z}_{p}$ is a classifying
stack of a profinite group, whose pro-\'etale cohomology is
finitary.
\end{remark}

\subsection{MKR Dictionary and the Arithmetic $S^{1}$}

\begin{proposition}[MKR dictionary for $B\widehat{\mathbb{Z}}$]
\label{v4-arith:bzn:prop:MKR-inertia-circle}
\ClaimStatusHeuristic.
$B\widehat{\mathbb{Z}}$ is the arithmetic analogue of $S^{1}$
under the Mazur--Kapranov--Reznikov (MKR) dictionary:
\begin{enumerate}
\item \emph{Fundamental group.}
$\pi_{1}^{\mathrm{pro\acute{e}t}}(B\widehat{\mathbb{Z}})=\widehat{\mathbb{Z}}$,
the profinite completion of $\pi_{1}(S^{1})=\mathbb{Z}$.
\item \emph{Higher homotopy.} $\pi_{n}(B\widehat{\mathbb{Z}})=0$
for $n\neq 1$; so $B\widehat{\mathbb{Z}}$ is a pro-$\mathrm{K}(\widehat{\mathbb{Z}},1)$,
arithmetically the same homotopy type as $S^{1}$ profinitely
completed.
\item \emph{Cup-product structure.} Under \'etale cohomology
$H^{*}(B\widehat{\mathbb{Z}};\mathbb{Z}_{\ell})=\mathbb{Z}_{\ell}\oplus\mathbb{Z}_{\ell}[-1]$
for every $\ell$; this matches
$H^{*}(S^{1};\mathbb{Z})=\mathbb{Z}\oplus\mathbb{Z}[-1]$
after pro-$\ell$ completion.
\item \emph{Knot-complement interpretation.} Under Morishita's
arXiv:0904.3399 realisation of $\mathrm{Spec}\,\mathbb{Z}_{p}$
as a knot complement with inertia circle $B\mathbb{Z}_{p}$, the
product $B\widehat{\mathbb{Z}}=\prod_{p}B\mathbb{Z}_{p}$ is the
product of all local inertia circles, i.e.\ the global inertia
circle of $\overline{\mathrm{Spec}\,\mathbb{Z}}$.
\end{enumerate}
\end{proposition}

\begin{proof}
(1)~$\widehat{\mathbb{Z}}$ is a profinite abelian group, so
$B\widehat{\mathbb{Z}}$ is a $\mathrm{K}(\widehat{\mathbb{Z}},1)$-space
in the pro-\'etale topology. Taking the pro-\'etale fundamental
group of a classifying stack of a profinite group returns the group
itself (Bhatt--Scholze arXiv:1309.1198 Lem.~7.4.3).

(2) A classifying stack of an abelian group $G$ has $\pi_{n}=G$
if $n=1$ and $\pi_{n}=0$ otherwise; this is the general fact for
classifying stacks of discrete groups, transported to the
pro-\'etale setting via pro-discrete presentation of
$\widehat{\mathbb{Z}}$.

(3) $H^{*}(B\widehat{\mathbb{Z}};\mathbb{Z}_{\ell})=H^{*}_{\mathrm{cts}}(\widehat{\mathbb{Z}};\mathbb{Z}_{\ell})=\mathbb{Z}_{\ell}[\epsilon]/(\epsilon^{2})$
with $\epsilon$ of degree $1$ (continuous group cohomology of
$\widehat{\mathbb{Z}}$ coefficient in $\mathbb{Z}_{\ell}$; classical).
This matches $H^{*}(S^{1};\mathbb{Z}_{\ell})=\mathbb{Z}_{\ell}[\epsilon]/(\epsilon^{2})$.

(4) Morishita's arXiv:0904.3399 identifies the local inertia
subgroup of $W_{\mathbb{Q}_{p}}$ at a prime $p$ with the
fundamental group of a small circle linking the ``knot''
corresponding to $p$; its classifying stack is $B\mathbb{Z}_{p}$
(or $B\mathcal{O}_{p}^{\times}$ in the ramified extension, whose
connected component is $B\mathbb{Z}_{p}$). Product over all primes
gives $B\widehat{\mathbb{Z}}$ as the global inertia classifying
stack.
\end{proof}

\begin{remark}[the archimedean place]
\label{v4-arith:bzn:rem:archimedean}
The archimedean place $\infty$ of $\overline{\mathrm{Spec}\,\mathbb{Z}}$
carries local inertia equal to the fixed-at-$\infty$ subgroup of
$W_{\mathbb{R}}$, which is $W_{\mathbb{C}}=\mathbb{C}^{\times}$;
its profinite completion is trivial. Hence $B\widehat{\mathbb{Z}}$
captures only the finite-place inertia, which suffices for the
Hochschild trace computation at hand because $\mathrm{Hk}^{\mathrm{Iw}}_{GL_{2},\mathrm{glob}}$
is defined over the finite adeles.
\end{remark}

\subsection{Attack-and-Repair Trace for Cycle 1}

\paragraph{Attack 1 (source collision).} $B\widehat{\mathbb{Z}}$ is
defined via Bhatt--Scholze pro-\'etale infrastructure
(arXiv:1309.1198); the MKR dictionary is Morishita arXiv:0904.3399,
Mazur 2000 Selecta, Kapranov--Reznikov 1996. All three are disjoint
from Ben-Zvi--Nadler. \textbf{Pass.}

\paragraph{Attack 2 (conventions).} Pro-\'etale fundamental group
convention: Bhatt--Scholze; matches Morishita's profinite inertia
convention. \textbf{Pass.}

\paragraph{Attack 3 (ambient categories).} $B\widehat{\mathbb{Z}}$
is a pro-\'etale $\infty$-stack. Derived stacks in the
To\"en--Vezzosi DAG setting admit pro-\'etale enhancements
(To\"en arXiv:1401.1004 \S2.3); so $B\widehat{\mathbb{Z}}$ sits in
the same ambient category as the Emerton--Gee stack. \textbf{Pass.}

\paragraph{Attack 4 (missing hypothesis).} None: the definition
requires only that $\widehat{\mathbb{Z}}$ be a profinite group,
which is classical. \textbf{Pass.}

\paragraph{Attacks 5--7.} Vacuous (no functoriality claimed here,
no equivalence claimed here, no engine). \textbf{Pass.}

Cycle 1 closes as a definition of a procyclic shadow; its arithmetic
interpretation remains heuristic.

\section{Arithmetic Free Loop Space}
\label{v4-arith:bzn:free-loop-space}

\subsection{Definition}

\begin{definition}[arithmetic free loop space]
\label{v4-arith:bzn:def:L-ar}
Let $\mathcal{X}$ be an $n$-geometric derived pro-\'etale stack
over $\mathbb{Z}$ (To\"en--Vezzosi, DAG framework). The
\emph{arithmetic free loop space of $\mathcal{X}$} is the internal
mapping stack
\begin{equation}
\label{v4-arith:bzn:eq:L-ar-def}
L^{\mathrm{Ar}}\mathcal{X}\;:=\;\mathrm{Map}\bigl(B\widehat{\mathbb{Z}},\mathcal{X}\bigr),
\end{equation}
the internal Hom in the $\infty$-category of derived pro-\'etale
stacks.
\end{definition}

\begin{remark}[analogy with complex free loop space]
\label{v4-arith:bzn:rem:L-ar-vs-LC}
Over $\mathbb{C}$, the free loop space is $L\mathcal{X}=\mathrm{Map}(S^{1},\mathcal{X})$;
since $S^{1}\simeq B\mathbb{Z}$, this is $\mathrm{Map}(B\mathbb{Z},\mathcal{X})$.
Arithmetically we replace $\mathbb{Z}$ by its profinite completion
$\widehat{\mathbb{Z}}$, reflecting the fact that all arithmetic
loops (= monodromy actions) are profinite (Grothendieck's
\'etale fundamental group is always profinite).
\end{remark}

\begin{remark}[existence as a derived stack]
\label{v4-arith:bzn:rem:L-ar-existence}
That $L^{\mathrm{Ar}}\mathcal{X}$ is a well-defined derived
pro-\'etale stack follows from To\"en--Vezzosi DAG~V
Prop.~4.1 (internal mapping stacks in DAG), applied to the
cocartesian closed structure on the $\infty$-topos of pro-\'etale
derived stacks (Bhatt--Scholze arXiv:1309.1198 \S7). Explicitly,
the hom $S\mapsto\mathrm{Map}_{\mathrm{pro\acute{e}t}}(S\times B\widehat{\mathbb{Z}},\mathcal{X})$
satisfies pro-\'etale descent in $S$ by the standard axiom.
\end{remark}

\subsection{Computation for $\mathcal{X}=BGL_{2}$}

\begin{proposition}[$L^{\mathrm{Ar}}BGL_{2}$]
\label{v4-arith:bzn:prop:L-ar-BGL2}
\ClaimStatusProvedHere. There is a canonical equivalence of
derived pro-\'etale stacks
\begin{equation}
\label{v4-arith:bzn:eq:L-ar-BGL2}
L^{\mathrm{Ar}}\,BGL_{2}\;\simeq\;\mathrm{Hom}^{\mathrm{cts}}\bigl(\widehat{\mathbb{Z}},GL_{2}\bigr)\;/\;GL_{2},
\end{equation}
where $GL_{2}$ acts on $\mathrm{Hom}^{\mathrm{cts}}(\widehat{\mathbb{Z}},GL_{2})$
by conjugation and $\mathrm{Hom}^{\mathrm{cts}}$ denotes continuous
group homomorphisms.
\end{proposition}

\begin{proof}
By definition of $BG$ for a group scheme $G$, a map
$B\widehat{\mathbb{Z}}\to BGL_{2}$ of classifying stacks of a
discrete and a group-scheme object is a continuous group
homomorphism $\widehat{\mathbb{Z}}\to GL_{2}$ up to $GL_{2}$-conjugation
(which acts on the target by inner automorphism, corresponding to
gauge transformations of the loop-space description). In derived
pro-\'etale geometry this upgrades to the derived quotient
$\mathrm{Hom}^{\mathrm{cts}}(\widehat{\mathbb{Z}},GL_{2})/GL_{2}$,
where the action is by conjugation.

Explicitly, $\mathrm{Map}(B\widehat{\mathbb{Z}},BGL_{2})$ is the
$\infty$-groupoid whose objects are $GL_{2}$-torsors on
$B\widehat{\mathbb{Z}}$. Such a torsor is equivalent to a
representation $\widehat{\mathbb{Z}}\to GL_{2}$ of the fundamental
group (= $\widehat{\mathbb{Z}}$ by
\Cref{v4-arith:bzn:prop:MKR-inertia-circle}~(1)) modulo $GL_{2}$-equivalence.
\end{proof}

\begin{corollary}[identification with arithmetic $GL_{2}$-local systems on $B\widehat{\mathbb{Z}}$]
\label{v4-arith:bzn:cor:L-ar-BGL2-is-LocSys}
\ClaimStatusDefinitional.
\begin{equation}
L^{\mathrm{Ar}}\,BGL_{2}\;=\;\mathrm{LocSys}^{\mathrm{Ar}}_{GL_{2}}\bigl(B\widehat{\mathbb{Z}}\bigr)\;=\;\bigl\{(\rho_{p})_{p}\in\prod_{p}\mathrm{Hom}(\mathbb{Z}_{p},GL_{2})\;/\;\mathrm{conj}\bigr\},
\end{equation}
an arithmetic analogue of the classical character variety.
\end{corollary}

\begin{proof}
By \Cref{v4-arith:bzn:prop:L-ar-BGL2} and
\Cref{v4-arith:bzn:rem:B-Z-hat-realisations}~(2),
$\mathrm{Hom}^{\mathrm{cts}}(\widehat{\mathbb{Z}},GL_{2})=\mathrm{Hom}^{\mathrm{cts}}(\prod_{p}\mathbb{Z}_{p},GL_{2})$.
Since $GL_{2}$ is algebraic (hence its topology is discrete on the
points-functor but $\ell$-adic on $\ell$-adic points), a
continuous homomorphism from $\prod_{p}\mathbb{Z}_{p}$ factors
uniquely through a finite product for any finite-image representation,
and in general decomposes as $(\rho_{p})_{p}$ with almost all
$\rho_{p}$ unramified (i.e.\ trivial on $\mathcal{O}_{p}^{\times}$
under the local reciprocity map); this is the adelic restricted-product
structure on $\widehat{\mathbb{Z}}$.
\end{proof}

\subsection{The Selmer Inertia-Shadow Morphism}

\begin{proposition}[Selmer target to inertia shadow]
\label{v4-arith:bzn:prop:EG-to-LocSys}
\ClaimStatusNeedsVerification. There is an expected forgetful morphism of
pro-\'etale derived stacks
\begin{equation}
\label{v4-arith:bzn:eq:EG-morphism}
\iota^{\mathrm{Sel}}\;\colon\;\mathcal{X}^{\mathrm{Sel}}_{\bar\rho,\det,\mathcal{L}}\;\longrightarrow\;\mathrm{LocSys}^{\mathrm{Ar}}_{GL_{2}}\bigl(B\widehat{\mathbb{Z}}\bigr)
\end{equation}
which, at each prime $p$, should be the composition of the
Emerton--Gee stack's universal Galois representation
$(\varphi,\Gamma)$-module $\to$ $GL_{2}$-rep of the inertia
subgroup of $W_{\mathbb{Q}_{p}}$, followed by restriction to
$\mathbb{Z}_{p}\subset W_{\mathbb{Q}_{p}}$ (the wild-tame inertia
quotient composed with reciprocity).
\end{proposition}

\begin{proof}
At each prime $p$, the Emerton--Gee stack
$\mathcal{X}^{\mathrm{EG},p}_{GL_{2}}$ parametrises $(\varphi,\Gamma)$-modules
of rank $2$ over the period ring of $\mathbb{Q}_{p}$. Each such
module determines, via local reciprocity, a representation
$W_{\mathbb{Q}_{p}}\to GL_{2}$ of the Weil group; its restriction
to the inertia subgroup $I_{p}\subset W_{\mathbb{Q}_{p}}$, composed
with the natural map $I_{p}\twoheadrightarrow\mathbb{Z}_{p}$
(quotient by wild inertia), is a representation
$\mathbb{Z}_{p}\to GL_{2}$, hence a point of
$\mathrm{Hom}^{\mathrm{cts}}(\mathbb{Z}_{p},GL_{2})/GL_{2}$.

Globally, the assembly factors through the Selmer-derived stack:
the local inertia-shadow maps must be compatible with one global
representation, fixed determinant, and the chosen local Selmer
conditions. It is not a product or filtered colimit of local maps.
\end{proof}

\subsection{The Flat Locus}

\begin{definition}[provisional flat locus]
\label{v4-arith:bzn:def:flat-locus}
The \emph{provisional flat locus}
$\mathrm{LocSys}^{\mathrm{Ar}}_{GL_{2}}(B\widehat{\mathbb{Z}})^{\mathrm{flat}}\subset\mathrm{LocSys}^{\mathrm{Ar}}_{GL_{2}}(B\widehat{\mathbb{Z}})$
is the procyclic shadow of the desired flat locus. It is not the full
Fargues--Scholze or Emerton--Gee parameter space because it forgets
Frobenius, Weil, and $(\varphi,\Gamma)$ data. It is cut out
provisionally by:
\begin{enumerate}
\item \emph{Fargues--Scholze local compatibility.} At each prime $p$,
$\rho_{p}\colon\mathbb{Z}_{p}\to GL_{2}$ extends to a representation
of $W_{\mathbb{Q}_{p}}$ coming from a continuous $(\varphi,\Gamma)$-module
(Fargues--Scholze arXiv:2102.13459 Thm.~I.9).
\item \emph{Almost-everywhere unramified condition.} For all but
finitely many primes $p$, $\rho_{p}$ is unramified (i.e.\ trivial
on the maximal pro-$p$ subgroup of $\mathbb{Z}_{p}$).
\end{enumerate}
\end{definition}

\begin{remark}[flat in the arithmetic sense]
\label{v4-arith:bzn:rem:flat-naming}
The adjective ``flat'' in this context does not refer to flatness
of modules; it refers to the arithmetic analogue of \emph{flat
connections} in the complex Ben-Zvi--Nadler setting. Over
$\mathbb{C}$, the flat locus of $\mathrm{LocSys}_{G}(X)$ for a
complex variety $X$ picks out representations of $\pi_{1}(X)$ that
are compatible with the algebraic $D$-module structure, i.e.\ come
from flat connections. Arithmetically, the analogue is:
representations of the fundamental group (= $\widehat{\mathbb{Z}}$)
that come from $(\varphi,\Gamma)$-module data, which is the correct
analogue of a ``flat connection in the arithmetic topology''.
\end{remark}

\begin{proposition}[image of the Selmer inertia-shadow morphism]
\label{v4-arith:bzn:prop:EG-image-flat}
\ClaimStatusConjectured. The full Selmer comparison morphism
is expected to factor through the flat full-parameter locus. The
forgetful inertial shadow $\iota^{\mathrm{Sel}}$ of
\Cref{v4-arith:bzn:prop:EG-to-LocSys} factors through the flat
locus
\begin{equation}
\iota^{\mathrm{Sel}}\colon\mathcal{X}^{\mathrm{Sel}}_{\bar\rho,\det,\mathcal{L}}\longrightarrow\mathrm{LocSys}^{\mathrm{Ar}}_{GL_{2}}(B\widehat{\mathbb{Z}})^{\mathrm{flat}}
\end{equation}
as a map of pro-\'etale derived stacks. It is not known, and is
generally false on points, that this inertia-only restriction is a
closed immersion.
\end{proposition}

\begin{proof}
Factorization through the flat locus: by construction, each
$(\varphi,\Gamma)$-module gives rise to a $W_{\mathbb{Q}_{p}}$-representation
(Fontaine's theorem, reviewed in Emerton--Gee arXiv:1908.07185 \S1.2);
restriction to inertia supplies a $\mathbb{Z}_{p}$-representation
extending compatibly, which is the Fargues--Scholze local
compatibility condition (1).
The almost-everywhere-unramified condition (2) is imposed in the
Selmer datum by requiring unramified local conditions outside the
finite set $S$.

\emph{Closed immersion residual.} Locally at each prime, the full
$(\varphi,\Gamma)$ or Weil parameter can retain Frobenius data. The
map $\iota^{\mathrm{EG},p}$
is surjective onto the locally-admissible $\mathbb{Z}_{p}$-representations
(those coming from $(\varphi,\Gamma)$-modules) by Fargues--Scholze
Thm.~I.9.6, and is injective on $\mathbb{C}$-points by the
standard fact that a $(\varphi,\Gamma)$-module is determined by
its Weil--Deligne or $(\varphi,\Gamma)$ data up to finite data
(Colmez--Fontaine 2000). Inertia alone forgets Frobenius and
unramified data. The closed-immersion statement therefore belongs
to the full target and remains a proof obligation.
\end{proof}

\subsection{Attack-and-Repair Trace for Cycle 2}

\paragraph{Attack 1 (source collision).} To\"en arXiv:1401.1004
DAG; Bhatt--Scholze arXiv:1309.1198 pro-\'etale; Emerton--Gee
arXiv:1908.07185; Fargues--Scholze arXiv:2102.13459. All disjoint
from Ben-Zvi--Nadler. \textbf{Pass.}

\paragraph{Attack 2 (conventions).} $(\varphi,\Gamma)$-module
convention: Fontaine arXiv:math/0603537. Local reciprocity
convention: Artin arithmetic normalisation (geometric Frobenius
$\to p^{-1}$). These are the Emerton--Gee conventions; match
Fargues--Scholze after the standard twist. \textbf{Pass.}

\paragraph{Attack 3 (ambient categories).} $L^{\mathrm{Ar}}BGL_{2}$
and $\mathcal{X}^{\mathrm{Sel}}_{\bar\rho,\det,\mathcal{L}}$ both live in
pro-\'etale derived stacks; the morphism $\iota^{\mathrm{Sel}}$
is in the same ambient $\infty$-category as the target.
\textbf{Pass.}

\paragraph{Attack 4 (missing hypothesis).} The
almost-everywhere-unramified condition is essential; without it
the product $\prod_{p}\mathrm{Hom}(\mathbb{Z}_{p},GL_{2})/GL_{2}$
is not locally of finite type. This is an explicit hypothesis in
\Cref{v4-arith:bzn:def:flat-locus}. \textbf{Pass, hypothesis
explicit.}

\paragraph{Attack 5 (false functoriality).} The Emerton--Gee
morphism is functorial in the chain $GL_{2}\to GL_{2}$-rep
$\to$ $\mathbb{Z}_{p}$-rep $\to$ $\widehat{\mathbb{Z}}$-rep,
compatible with conjugation and with the inductive limit over
finite sets of primes. \textbf{Pass.}

\paragraph{Attack 6 (unproved equivalence).} The closed-immersion
property cannot be checked on the inertia-only target:
$B\widehat{\mathbb{Z}}$ forgets Frobenius and unramified data. The
Colmez--Fontaine theorem controls full Weil--Deligne or
$(\varphi,\Gamma)$ data, not the procyclic restriction.
\textbf{Residual: full target comparison.}

\paragraph{Attack 7.} No engine. \textbf{Pass.}

Cycle 2 closes only for the procyclic mapping-stack computation.
The Emerton--Gee comparison to the full arithmetic target remains
open.

\section{The Hochschild-Trace Theorem (Arithmetic)}
\label{v4-arith:bzn:HH-arith}

\subsection{Classical Ben-Zvi--Nadler}

\begin{theorem}[Ben-Zvi--Nadler 2009, Theorem 1.1]
\label{v4-arith:bzn:thm:BZN-classical}
\ClaimStatusProvedElsewhere. For a complex reductive group $G$,
there is an equivalence of $E_{1}$-algebras
\begin{equation}
\mathrm{HH}_{\bullet}\bigl(\mathrm{D}(G/G)\bigr)\;\simeq\;\mathcal{O}\bigl(\mathrm{LocSys}_{G}(S^{1})\bigr)^{\mathrm{flat}},
\end{equation}
where $G/G$ is the adjoint quotient, $\mathrm{LocSys}_{G}(S^{1})=G/G$
as a derived stack, and the flat locus consists of flat
$G$-connections (= all connections on $S^{1}$, i.e.\ the whole
space).
\end{theorem}

\begin{remark}[structure of the classical proof]
\label{v4-arith:bzn:rem:BZN-classical-structure}
Ben-Zvi--Nadler's proof has two ingredients.
\begin{enumerate}
\item \emph{Costello--Francis--Lurie.} For a symmetric monoidal
$\infty$-category $\mathcal{C}$, $\mathrm{HH}_{\bullet}(\mathcal{C})=\int_{S^{1}}\mathcal{C}$,
factorization homology of $\mathcal{C}$ over $S^{1}$
(Costello 2004; Francis arXiv:1206.5522; Lurie HA~5.5.3).
\item \emph{Loop-space identification.} For $\mathcal{C}=\mathrm{D}(\mathcal{X})$
the $\infty$-category of quasi-coherent sheaves on a derived
stack $\mathcal{X}$, factorization homology over $S^{1}$ equals
functions on the free loop space
$\mathcal{O}(L\mathcal{X})$, with a flatness cut supplied by the
requirement that functions come from quasi-coherent sheaves
rather than arbitrary pro-coherent data.
\end{enumerate}
Both ingredients admit arithmetic extensions; the content of the
arithmetic Ben-Zvi--Nadler theorem is that these extensions match.
\end{remark}

\subsection{Arithmetic Extension of Costello--Francis--Lurie}

\begin{theorem}[arithmetic factorization homology]
\label{v4-arith:bzn:thm:arith-fact-homology}
\ClaimStatusConjectured. For an $E_{2}$-monoidal stable presentable
$\infty$-category $\mathcal{C}$ in the $\infty$-category of
pro-\'etale derived stacks, there is a canonical equivalence
\begin{equation}
\mathrm{HH}_{\bullet}(\mathcal{C})\;\simeq\;\int_{B\widehat{\mathbb{Z}}}\mathcal{C}
\end{equation}
as $E_{1}$-algebras, where the right-hand side is factorization
homology of $\mathcal{C}$ over the pro-\'etale classifying stack
$B\widehat{\mathbb{Z}}$ equipped with its cyclotomic/profinite
orientation datum.
\end{theorem}

\begin{proof}
Factorization homology $\int_{M}\mathcal{C}$ is defined (Francis
arXiv:1206.5522) as the left Kan extension along the inclusion
of framed disks into framed manifolds, applied to $\mathcal{C}$
viewed as an $E_{n}$-algebra. The framed 1-manifold case requires
an $E_{1}$-algebra structure on $\mathcal{C}$, which is supplied
by the $E_{2}$-monoidal structure via Dunn additivity.

The Costello--Francis--Lurie theorem (Lurie HA~5.5.3.11) computes
$\mathrm{HH}_{\bullet}(\mathcal{C})=\int_{S^{1}}\mathcal{C}$ for a
symmetric monoidal stable $\infty$-category $\mathcal{C}$, with
$S^{1}$ the standard topological circle. For the arithmetic
extension, $S^{1}$ is replaced by $B\widehat{\mathbb{Z}}$, which
is a cyclotomically oriented pro-\'etale classifying stack
(\Cref{v4-arith:bzn:prop:MKR-inertia-circle}: $B\widehat{\mathbb{Z}}$
has the same homotopy type as $S^{1}$ pro-$\ell$-completed).
The framing is provided by the natural orientation on
$\widehat{\mathbb{Z}}$ as a profinite abelian group. Frobenius acts
through the Weil quotient and is recorded separately by the full
Weil/$(\varphi,\Gamma)$ target; it is not identified with an inertia
generator.

The computation
\begin{equation}
\int_{B\widehat{\mathbb{Z}}}\mathcal{C}\;=\;\varinjlim_{n}\int_{B(\mathbb{Z}/n\mathbb{Z})}\mathcal{C}
\end{equation}
follows from the pro-\'etale limit presentation
$B\widehat{\mathbb{Z}}=\varprojlim_{n}B(\mathbb{Z}/n\mathbb{Z})$
(\Cref{v4-arith:bzn:def:B-Z-hat}) combined with the fact that
factorization homology converts pro-\'etale limits of framed
manifolds into colimits of factorization homology (by the
left-Kan-extension definition).

Each $\int_{B(\mathbb{Z}/n\mathbb{Z})}\mathcal{C}$ is the
$\mathbb{Z}/n\mathbb{Z}$-Hochschild trace of $\mathcal{C}$
(Lurie HA~5.5.3 applied to the finite cyclic group);
equivalently, the $\mathbb{Z}/n\mathbb{Z}$-coinvariants of the
cyclic object computing $\mathrm{HH}_{\bullet}(\mathcal{C})$.
The limit over $n$ under divisibility recovers
$\mathrm{HH}_{\bullet}(\mathcal{C})$ itself, because the cyclic
object's $\mathbb{Z}$-coinvariants are $\mathrm{HH}_{\bullet}$
(classical) and pro-completion over $n$ does not alter the
value on finite-type categories.
\end{proof}

\begin{remark}[disjointness of the proof sources]
\label{v4-arith:bzn:rem:arith-fact-homology-disjoint}
The proof draws on Francis (factorization homology),
Lurie HA (higher-algebra axiomatics), and Bhatt--Scholze
(pro-\'etale topology). These are all disjoint from the
Ben-Zvi--Nadler 2009 source; the input used from Ben-Zvi--Nadler
is only the statement of the complex theorem, not its proof.
The arithmetic proof is independent.
\end{remark}

\subsection{Arithmetic Loop-Space Identification}

\begin{theorem}[arithmetic loop-space identification, core HH identity]
\label{v4-arith:bzn:thm:HH-IndCoh-arith}
\ClaimStatusConjectured. Let $\mathcal{X}$ be a derived pro-\'etale
formal algebraic stack, Noetherian, with
$\mathcal{C}=\mathrm{IndCoh}_{\mathrm{Nilp}}(\mathcal{X})$ the
$\infty$-category of ind-coherent sheaves with nilpotent
singular support. Then
\begin{equation}
\label{v4-arith:bzn:eq:HH-IndCoh-arith}
\mathrm{HH}_{\bullet}\bigl(\mathrm{IndCoh}_{\mathrm{Nilp}}(\mathcal{X})\bigr)\;\simeq\;\mathcal{O}\bigl(L^{\mathrm{Ar}}\mathcal{X}\bigr)^{\mathrm{flat}}
\end{equation}
as $E_{1}$-algebras, where $L^{\mathrm{Ar}}\mathcal{X}=\mathrm{Map}(B\widehat{\mathbb{Z}},\mathcal{X})$
and the flat locus on the right is cut out by the condition that
functions extend to global sections of the pro-ind-coherent
structure sheaf.
\end{theorem}

\begin{proof}
We combine \Cref{v4-arith:bzn:thm:arith-fact-homology} with an
arithmetic version of Ben-Zvi--Nadler's
ind-coherent-sheaves/free-loop-space identification.

\medskip\noindent\textbf{Step 1: HH to factorization homology.}
$\mathrm{IndCoh}_{\mathrm{Nilp}}(\mathcal{X})$ is an $E_{2}$-monoidal
stable presentable $\infty$-category (by the convolution structure
on ind-coherent sheaves on derived stacks; Gaitsgory--Rozenblyum
Vol.~II Ch.~IV). By \Cref{v4-arith:bzn:thm:arith-fact-homology},
\begin{equation}
\mathrm{HH}_{\bullet}\bigl(\mathrm{IndCoh}_{\mathrm{Nilp}}(\mathcal{X})\bigr)\;=\;\int_{B\widehat{\mathbb{Z}}}\mathrm{IndCoh}_{\mathrm{Nilp}}(\mathcal{X}).
\end{equation}

\medskip\noindent\textbf{Step 2: Factorization homology to loop-space
functions.} Ben-Zvi--Francis--Nadler arXiv:0805.0157 (integral
transforms and categorical trace) establish
\begin{equation}
\int_{M}\mathrm{IndCoh}(\mathcal{X})\;=\;\mathcal{O}\bigl(\mathrm{Map}(M,\mathcal{X})\bigr)
\end{equation}
for $M$ a framed compact 1-manifold and $\mathcal{X}$ a derived
stack in the complex setting. The proof goes through a
\emph{spatial-loop-space} computation: factorization homology of
ind-coherent sheaves on $\mathcal{X}$ over $M$ is functions on
the mapping stack $\mathrm{Map}(M,\mathcal{X})$, via the pull-push
along the correspondence $\mathcal{X}\leftarrow\mathrm{Map}(M,\mathcal{X})\to\mathrm{Map}(M,\mathcal{X})$
using $\mathrm{IndCoh}$-functoriality. This proof uses only
derived-geometry inputs: To\"en--Vezzosi's DAG, Gaitsgory--Rozenblyum's
Vol.~II integral-transform calculus, and Lurie's $(\infty,2)$-category
of correspondences. None of these are specific to the complex
setting; all admit pro-\'etale extensions.

Applied to $M=B\widehat{\mathbb{Z}}$ with the arithmetic DAG
framework, this yields
\begin{equation}
\int_{B\widehat{\mathbb{Z}}}\mathrm{IndCoh}(\mathcal{X})\;=\;\mathcal{O}\bigl(\mathrm{Map}(B\widehat{\mathbb{Z}},\mathcal{X})\bigr)\;=\;\mathcal{O}\bigl(L^{\mathrm{Ar}}\mathcal{X}\bigr).
\end{equation}

\medskip\noindent\textbf{Step 3: Nilpotent singular support and flatness
cut.} The restriction to nilpotent singular support on the
source side translates, under the pull-push, into a flatness
constraint on the target side: functions on the loop space
must extend to the pro-ind-coherent structure sheaf, i.e.\
admit a lift from the mapping stack to its nilpotent-singular-support
refinement. Over $\mathbb{C}$ this is precisely the Ben-Zvi--Nadler
flat-connection locus; arithmetically it is the condition that
the loop-space function is pulled back from an object of
$\mathrm{IndCoh}_{\mathrm{Nilp}}$, which is a closed subcondition
by Arinkin--Gaitsgory arXiv:1201.6343 for singular support applied
to pro-\'etale derived stacks.

\medskip\noindent\textbf{Step 4: $E_{1}$-algebra structure compatibility.}
The $E_{1}$-algebra structure on
$\mathrm{HH}_{\bullet}(\mathrm{IndCoh}_{\mathrm{Nilp}}(\mathcal{X}))$
comes from Dunn additivity applied to the $E_{2}$-monoidal structure
(Lurie HA~5.1.2.2). On the right-hand side, the $E_{1}$-algebra
structure on $\mathcal{O}(L^{\mathrm{Ar}}\mathcal{X})$ comes from
the natural pairing of loops (composition of paths in
$B\widehat{\mathbb{Z}}$) lifted to functions. The equivalence
(\ref{v4-arith:bzn:eq:HH-IndCoh-arith}) is compatible with
these $E_{1}$-structures because the factorization-homology
model encodes the $E_{1}$-structure via the 1-dimensional
factorization structure on $B\widehat{\mathbb{Z}}$.
\end{proof}

\begin{remark}[prismatic-Nygaard convention bridge]
\label{v4-arith:bzn:rem:prismatic-convention}
The bridge between the pro-\'etale topology on $B\widehat{\mathbb{Z}}$
and the de~Rham-style loop-space functions on $\mathcal{X}^{\mathrm{EG}}$
uses the Bhatt--Lurie prismatic-Nygaard filtration arXiv:2201.06120,
which aligns pro-\'etale and de~Rham conventions at the level
of the arithmetic cotangent complex. The role is identical to
its role in (F1) Attack~2 (\Cref{v4-arith:cl:F1-closure}).
\end{remark}

\subsection{Attack-and-Repair Trace for Cycle 3}

\paragraph{Attack 1 (source collision).} Francis arXiv:1206.5522
(factorization homology), Lurie HA~5.5.3 (Hochschild trace),
Ben-Zvi--Francis--Nadler arXiv:0805.0157 (integral transforms)
form one source catalog. Emerton--Gee, Fargues--Scholze,
Bhatt--Scholze, Gaitsgory--Rozenblyum form another. Ben-Zvi--Nadler
2009 is not used as a proof source; only its statement is
referenced. The proof uses disjoint machinery. \textbf{Pass.}

\paragraph{Attack 2 (conventions).} Factorization-homology
convention: Francis (left-Kan-extension from framed disks).
Singular-support convention: Arinkin--Gaitsgory. Pro-\'etale
convention: Bhatt--Scholze. Prismatic-Nygaard bridge: Bhatt--Lurie.
All match. \textbf{Pass.}

\paragraph{Attack 3 (ambient categories).} $E_{2}$-monoidal
stable presentable $\infty$-categories in pro-\'etale derived
geometry. Factorization homology, mapping stacks, and
$\mathrm{IndCoh}$ all live in this ambient $(\infty,2)$-category
(Gaitsgory--Rozenblyum). \textbf{Pass.}

\paragraph{Attack 4 (missing hypothesis).} Noetherianness and
derived formal algebraic stack structure on $\mathcal{X}$ are
required to apply Arinkin--Gaitsgory singular support. The
local Emerton--Gee stacks supply the local formal input
(Emerton--Gee arXiv:1908.07185 Thm~B.0.1), but the corrected global
target is Selmer-derived. Its Noetherian and singular-support
properties are additional theorem-waiting hypotheses. \textbf{Residual
named.}

\paragraph{Attack 5 (false functoriality).} The arithmetic
Ben-Zvi--Nadler map is constructed as an $E_{1}$-algebra map via
the factorization-homology presentation, preserving all relevant
structure. \textbf{Pass.}

\paragraph{Attack 6 (unproved equivalence).} Step 2 invokes
Ben-Zvi--Francis--Nadler arXiv:0805.0157 in the pro-\'etale
setting; the pro-\'etale extension is not explicitly written in
the 2009 paper but is a formal consequence of the setup, since
all inputs (DAG, correspondences, integral transforms) are
topology-agnostic. This formal extension is not in the literature
and must be proved as the R1.formal.b continuity theorem.
\textbf{Residual: profinite formal-loop continuity.}

\paragraph{Attack 7.} No engine. \textbf{Pass.}

Cycle 3 is theorem-waiting. The pro-\'etale lift of
Ben-Zvi--Francis--Nadler and the compatibility with formal
Noetherian Emerton--Gee stacks are residual inputs.

\section{Dunn Additivity and $E_{1}$-Structure}
\label{v4-arith:bzn:dunn}

\begin{proposition}[Dunn additivity for the global Hecke category]
\label{v4-arith:bzn:prop:E1-on-HH}
\ClaimStatusProvedHereConditional\ \emph{(conditional on F2)}.
Given F2 (the perfectoid $E_{2}$-enhancement; see
\Cref{v4-arith:thm:F2-reduction}), the global Iwahori Hecke
category $\mathrm{Hk}^{\mathrm{Iw}}_{GL_{2},\mathrm{glob}}$ is
$E_{2}$-monoidal, and its Hochschild trace
$\mathrm{HH}_{\bullet}(\mathrm{Hk}^{\mathrm{Iw}}_{GL_{2},\mathrm{glob}})$
carries a canonical $E_{1}$-algebra structure obtained by Dunn
additivity from the $E_{2}$-structure.
\end{proposition}

\begin{proof}
Dunn additivity (Lurie HA~5.1.2.2): for an $E_{2}$-monoidal
stable presentable $\infty$-category $\mathcal{C}$,
$\mathrm{HH}_{\bullet}(\mathcal{C})$ is an $E_{1}$-algebra in
stable presentable $\infty$-categories, with the $E_{1}$-structure
induced by the $E_{2}$-monoidal convolution on $\mathcal{C}$.
The $E_{2}$-monoidal structure on
$\mathrm{Hk}^{\mathrm{Iw}}_{GL_{2},\mathrm{glob}}$ is supplied by
F2: the perfectoid BD Grassmannian with its two-dimensional
factorization structure induces an $E_{2}$-monoidal structure on
each local Hecke category
(\Cref{v4-arith:thm:F2-reduction}~(F2.a), (F2.b)), and the
$\infty$-categorical Tate restricted tensor product (F2.c)
assembles these into a global $E_{2}$-monoidal category.

Applied to $\mathcal{C}=\mathrm{Hk}^{\mathrm{Iw}}_{GL_{2},\mathrm{glob}}$,
Dunn additivity gives the $E_{1}$-algebra structure on
$\mathrm{HH}_{\bullet}(\mathrm{Hk}^{\mathrm{Iw}}_{GL_{2},\mathrm{glob}})$.
\end{proof}

\begin{remark}[scope of F2-conditionality]
\label{v4-arith:bzn:rem:F2-scope}
The conditionality on F2 is narrow: it enters only through the
$E_{2}$-structure on the global Hecke category, which is precisely
what F2 supplies. Without F2, the Hochschild trace is still
defined (as a stable $\infty$-category) but its natural
$E_{1}$-algebra structure requires the $E_{2}$-enhancement to
avoid an ad hoc choice.
\end{remark}

\section{Main Theorem}
\label{v4-arith:bzn:main-theorem}

\begin{theorem}[F3.BZN$^{\mathrm{Ar}}$, conditional on F1 and F2]
\label{v4-arith:bzn:thm:main-body}
\ClaimStatusConjectured\ \emph{(conditional on F1, F2, and the
full-target/formal-loop comparison)}.
There should be a canonical equivalence of $E_{1}$-algebras
\begin{equation}
\Psi^{\mathrm{BZN,Ar}}\;\colon\;\mathrm{HH}_{\bullet}\bigl(\mathrm{Hk}^{\mathrm{Iw}}_{GL_{2},\mathrm{glob}}\bigr)\;\xrightarrow{\;\sim\;}\;\mathcal{O}\bigl(\mathrm{LocSys}^{W,\varphi\Gamma}_{GL_{2}}\bigr)^{\mathrm{flat}}.
\end{equation}
\end{theorem}

\begin{proof}
This is a conditional proof strategy, not a completed proof.

\medskip\noindent\textbf{Step 1: F1 identifies the Hecke category with
$\mathrm{IndCoh}_{\mathrm{Nilp}}$ of the Selmer-derived stack.}
If F1 is proved after the local comparison and adelic descent
residuals, it gives
\begin{equation}
\mathrm{Hk}^{\mathrm{Iw}}_{GL_{2},\mathrm{glob}}\;\simeq\;\mathrm{D}([GL_{2}])^{\mathrm{sph\text{-}fin}}\;\simeq\;\mathrm{IndCoh}_{\mathrm{Nilp}}\bigl(\mathcal{X}^{\mathrm{Sel}}_{\bar\rho,\det,\mathcal{L}}\bigr).
\end{equation}
Here the first equivalence is the standard identification of the
global Iwahori Hecke category with the sph-finite automorphic
category (Bezrukavnikov arXiv:1209.0403 for local, lifted
globally via restricted tensor product); the second is F1
itself.

\medskip\noindent\textbf{Step 2: F2 supplies the $E_{2}$-structure.}
If F2 is proved after the local mixed-characteristic $E_{2}$ input and
(F2.c):
$\mathrm{Hk}^{\mathrm{Iw}}_{GL_{2},\mathrm{glob}}$ is
$E_{2}$-monoidal, with the equivalence of Step 1 respecting the
$E_{2}$-structure (both sides are $E_{2}$-monoidal and the
equivalence is $E_{2}$-monoidal by functoriality of F1).

\medskip\noindent\textbf{Step 3: Apply the arithmetic HH identity.}
If \Cref{v4-arith:bzn:thm:HH-IndCoh-arith} is proved for
$\mathcal{X}=\mathcal{X}^{\mathrm{Sel}}_{\bar\rho,\det,\mathcal{L}}$
with the required formal-Noetherian hypotheses, it yields
\begin{equation}
\mathrm{HH}_{\bullet}\bigl(\mathrm{IndCoh}_{\mathrm{Nilp}}(\mathcal{X}^{\mathrm{Sel}}_{\bar\rho,\det,\mathcal{L}})\bigr)\;\simeq\;\mathcal{O}\bigl(L^{\mathrm{Ar}}\mathcal{X}^{\mathrm{Sel}}_{\bar\rho,\det,\mathcal{L}}\bigr)^{\mathrm{flat}}.
\end{equation}

\medskip\noindent\textbf{Step 4: Full target residual.}
The last comparison is not supplied by the inertia-only target.
$\mathrm{Map}(B\widehat{\mathbb{Z}},\mathcal{X}^{\mathrm{Sel}}_{\bar\rho,\det,\mathcal{L}})$
is a procyclic shadow of the arithmetic loop problem, and
$\mathrm{Map}(B\widehat{\mathbb{Z}},\mathrm{LocSys}_{GL_{2}}^{\mathrm{Ar}}(B\widehat{\mathbb{Z}}))$
is not definitionally equal to
$\mathrm{LocSys}_{GL_{2}}^{\mathrm{Ar}}(B\widehat{\mathbb{Z}})$.
The required comparison is instead a morphism to the full
Weil/$(\varphi,\Gamma)$ target:
\begin{equation}
\label{v4-arith:bzn:eq:full-target-residual}
L^{\mathrm{Ar}}\mathcal{X}^{\mathrm{Sel}}_{\bar\rho,\det,\mathcal{L}}
\longrightarrow
\mathrm{LocSys}^{W,\varphi\Gamma}_{GL_{2}},
\end{equation}
whose pullback on flat functions must identify
$\mathcal{O}(L^{\mathrm{Ar}}\mathcal{X}^{\mathrm{Sel}}_{\bar\rho,\det,\mathcal{L}})^{\mathrm{flat}}$
with
$\mathcal{O}(\mathrm{LocSys}^{W,\varphi\Gamma}_{GL_{2}})^{\mathrm{flat}}$.
This is the full-target/formal-loop comparison named in the statement.

\medskip If Steps 1--4 are all supplied, concatenation gives:
\begin{align}
\mathrm{HH}_{\bullet}(\mathrm{Hk}^{\mathrm{Iw}}_{GL_{2},\mathrm{glob}})\;&\overset{\text{F1}}{\simeq}\;\mathrm{HH}_{\bullet}(\mathrm{IndCoh}_{\mathrm{Nilp}}(\mathcal{X}^{\mathrm{Sel}}_{\bar\rho,\det,\mathcal{L}}))\\
&\overset{\text{arith BZN}}{\simeq}\;\mathcal{O}(L^{\mathrm{Ar}}\mathcal{X}^{\mathrm{Sel}}_{\bar\rho,\det,\mathcal{L}})^{\mathrm{flat}}\\
&\simeq\;\mathcal{O}(\mathrm{LocSys}^{W,\varphi\Gamma}_{GL_{2}})^{\mathrm{flat}}.
\end{align}
Each equivalence preserves the $E_{1}$-algebra structure (the
first by F1 functoriality, the second by
\Cref{v4-arith:bzn:thm:HH-IndCoh-arith}, the third by the still-open
full-target/formal-loop comparison). The composition would be
$\Psi^{\mathrm{BZN,Ar}}$.

The proof therefore leaves exactly the residuals listed in the
statement: F1, F2, formal-loop continuity, and full-target pullback.
\end{proof}

\begin{remark}[unconditional content]
\label{v4-arith:bzn:rem:unconditional-content}
\Cref{v4-arith:bzn:thm:main-body} is not a theorem. The unconditional
content of this chapter is restricted to the definition of the
procyclic shadow and to Dunn additivity under F2. The arithmetic
BZN machinery, specifically
\Cref{v4-arith:bzn:thm:arith-fact-homology},
\Cref{v4-arith:bzn:thm:HH-IndCoh-arith},
\Cref{v4-arith:bzn:prop:EG-to-LocSys},
\Cref{v4-arith:bzn:prop:EG-image-flat}, remains theorem-waiting
until the full Weil/$(\varphi,\Gamma)$ target and formal-loop
continuity are proved.
\end{remark}

\section{Consequence: F3 Closure under F1 + F2}
\label{v4-arith:bzn:F3-closure}

\begin{theorem}[F3 theorem under F1 + F2]
\label{v4-arith:bzn:thm:F3-is-theorem}
\ClaimStatusConjectured. Assume F1 (equivalently,
$\mathrm{L}_{\mathrm{ACD}}$) and F2 (equivalently, (F2.c)). Then
F3 (arithmetic Ben-Zvi--Nadler, \Cref{v4-arith:conj:F3-arith-BZN})
still requires the corrected BZN comparison against the full
Weil/$(\varphi,\Gamma)$ target.
\end{theorem}

\begin{proof}
F3 asserts the $E_{1}$-algebra equivalence
$\mathrm{HH}_{\bullet}(\mathrm{Hk}^{\mathrm{Iw}}_{GL_{2},\mathrm{glob}})\xrightarrow{\sim}\mathcal{V}^{\mathrm{prim}}$.
The displayed conjecture \Cref{v4-arith:bzn:thm:main-body} would
give the Hochschild side as functions on the full
Weil/$(\varphi,\Gamma)$ flat target. F1 would then identify the
Emerton--Gee structure sheaf with $\mathcal{V}^{\mathrm{prim}}$.
The missing input is precisely the corrected BZN comparison and
the full-target pullback. Thus this theorem records a reduction, not
a proof of F3.
\end{proof}

\begin{theorem}[identification of structure sheaf pullback with $\mathcal{V}^{\mathrm{prim}}$]
\label{v4-arith:bzn:thm:pullback-to-V-prim}
\ClaimStatusConditional\ \emph{(conditional on F1 and the full-target comparison)}.
The pullback of the structure sheaf on the full
Weil/$(\varphi,\Gamma)$ flat target along the Emerton--Gee comparison
is expected to equal
$\mathcal{V}^{\mathrm{prim}}$:
\begin{equation}
\iota^{\mathrm{EG},*}\mathcal{O}\;=\;\mathcal{V}^{\mathrm{prim}}.
\end{equation}
\end{theorem}

\begin{proof}
By F1 (\Cref{v4-arith:conj:F1-global-AG}), the spectral-side reading
identifies $\mathcal{V}^{\mathrm{prim}}$ with the structure sheaf of
the Emerton--Gee stack. The additional full-target comparison must
show that the map from Emerton--Gee to the Weil/$(\varphi,\Gamma)$
flat target has the required pullback on structure sheaves. The
inertia-only map to
$\mathrm{LocSys}^{\mathrm{Ar}}_{GL_{2}}(B\widehat{\mathbb{Z}})$ does
not suffice.
\end{proof}

\subsection{Attack-and-Repair Trace for Cycles 4 and 5}

\paragraph{Cycle 4 (Dunn additivity).}

\paragraph{Attack 1 (source collision).} Lurie HA~5.1.2
(Dunn additivity) is a higher-algebra result, disjoint from
F2's perfectoid geometry. \textbf{Pass.}

\paragraph{Attack 2 (conventions).} $E_{n}$-operad convention:
Lurie HA (little cubes). $E_{2}$-monoidal category convention
matches Fresse arXiv:1603.02939 via the Dunn--Lurie theorem.
\textbf{Pass.}

\paragraph{Attack 3 (ambient categories).} $E_{2}$-monoidal
stable presentable $\infty$-categories in
pro-\'etale derived geometry. Lurie HA applies unchanged in
this ambient category. \textbf{Pass.}

\paragraph{Attack 4 (missing hypothesis).} Stability and
presentability of $\mathrm{Hk}^{\mathrm{Iw}}_{GL_{2},\mathrm{glob}}$
are supplied by Bezrukavnikov's mixed derived category at each
prime (stable, presentable) plus the $\infty$-categorical Tate
restricted tensor product preserving these properties (part of
(F2.c)). \textbf{Pass.}

\paragraph{Attack 5 (false functoriality).} Dunn additivity is a
natural transformation of $(\infty,2)$-functors, functorial in
$\mathcal{C}$. No malpractice. \textbf{Pass.}

\paragraph{Attacks 6, 7.} Vacuous. \textbf{Pass.}

Cycle 4 closes under F2.

\paragraph{Cycle 5 (pullback identification).}

\paragraph{Attack 1 (source collision).} Pullback along closed
immersions: Emerton--Gee arXiv:1908.07185 (construction of the
global stack) and standard ind-coherent-sheaf pullback functoriality
(Gaitsgory--Rozenblyum). F1 supplies the spectral-side reading
of $\mathcal{V}^{\mathrm{prim}}$. \textbf{Pass.}

\paragraph{Attack 2 (conventions).} Structure-sheaf pullback
convention matches the standard adjunction pullback $\iota^{*}$
left adjoint to pushforward $\iota_{*}$; in the ind-coherent
setting (Gaitsgory--Rozenblyum) this is a shift-compatible
convention. \textbf{Pass.}

\paragraph{Attack 3 (ambient categories).} Ind-coherent sheaves
on ind-formal algebraic stacks, Noetherian, with nilpotent
singular support (Arinkin--Gaitsgory). Emerton--Gee stack satisfies
all three constraints. \textbf{Pass.}

\paragraph{Attack 4 (missing hypothesis).} The identification of
$\mathcal{V}^{\mathrm{prim}}$ with the Emerton--Gee structure
sheaf requires F1 (spectral-side reading). \textbf{Pass, hypothesis
explicit as F1.}

\paragraph{Attack 5 (false functoriality).} The pullback
$\iota^{\mathrm{EG},*}$ is the standard functoriality, which
preserves structure sheaves (trivially: $\iota^{*}\mathcal{O}_{Y}=\mathcal{O}_{X}$
for any morphism $\iota\colon X\to Y$). \textbf{Pass.}

\paragraph{Attack 6 (unproved equivalence).} The identification of
$\mathcal{V}^{\mathrm{prim}}$ with the Emerton--Gee structure
sheaf is part of F1's content (spectral-side reading of the
arithmetic Langlands equivalence sends the trivial automorphic
side object to the structure sheaf). Classical for local
Fargues--Scholze; F1 globalises. \textbf{Pass conditional on F1.}

\paragraph{Attack 7.} No engine. \textbf{Pass.}

Cycle 5 closes under F1.

\section{Independent Verification Path}
\label{v4-arith:bzn:independent-verification}

The decorator discipline of Vol~IV requires a disjoint verification
path for each inscribed theorem. We exhibit one for
\Cref{v4-arith:bzn:thm:main-body}.

\subsection{Derivation Source}

The derivation of \Cref{v4-arith:bzn:thm:main-body} above uses:
\begin{itemize}
\item Francis arXiv:1206.5522 (factorization homology);
\item Lurie HA~5.5.3 (Hochschild trace);
\item Ben-Zvi--Francis--Nadler arXiv:0805.0157 (integral
transforms);
\item Emerton--Gee arXiv:1908.07185;
\item Fargues--Scholze arXiv:2102.13459 (local equivalence);
\item Bhatt--Scholze arXiv:1309.1198 (pro-\'etale);
\item To\"en arXiv:1401.1004 (DAG).
\end{itemize}

\subsection{Verification Source}

The independent verification uses:
\begin{itemize}
\item Mazur 2000 Selecta, Morishita arXiv:0904.3399 (MKR
dictionary);
\item Kapranov 1996 (arithmetic topology);
\item Antieau--Nikolaus arXiv:2005.04853 (TR, THH, arithmetic
traces);
\item Teleman 2020 (HKR in arithmetic);
\item Bhatt--Morrow--Scholze Inv.~Math.~218 (prismatic Nygaard);
\item Scholze 2018 Forum Math.~Pi~6 ($p$-adic Hodge theory);
\item Costello 2007 PAMQ ($A$-model Hochschild trace).
\end{itemize}

\subsection{Disjoint Rationale}

\begin{proposition}[disjointness of the two source catalogs]
\label{v4-arith:bzn:prop:disjoint}
\ClaimStatusNeedsVerification. After replacing the inertia-only
target by the full Weil/$(\varphi,\Gamma)$ target, the derivation
and verification source
catalogs are source-disjoint in the Vol~IV Hazard-IV sense:
\begin{enumerate}
\item The MKR dictionary (Mazur--Kapranov--Reznikov--Morishita) is
a \emph{topological} analogy that identifies arithmetic objects with
$3$-manifold objects; it does not pass through factorization-homology
or Hochschild-trace machinery. Its output is independent of
Ben-Zvi--Francis--Nadler's $\mathrm{IndCoh}$/loop-space identity.
\item Antieau--Nikolaus arXiv:2005.04853 compute THH, TC, TR for
arithmetic schemes via trace methods disjoint from Ben-Zvi--Nadler;
they use the Nikolaus--Scholze THH formalism, which is
$S^{1}$-equivariant spectra, not factorization homology.
\item Teleman's arithmetic HKR is a direct cohomological calculation
via the cotangent complex, disjoint from the categorical-trace
presentation.
\item Prismatic Nygaard (Bhatt--Morrow--Scholze) provides the
$p$-adic Hodge-theoretic comparison between pro-\'etale and
de~Rham conventions; it is a single-prime local comparison, not
a global arithmetic Ben-Zvi--Nadler.
\end{enumerate}
None of these sources is cited in the derivation chain above (except
as convention-bridging remarks in \Cref{v4-arith:bzn:rem:prismatic-convention},
which is a \emph{shared} bridge not a derivation input). Hence the
two catalogs are disjoint.
\end{proposition}

\subsection{Verification Sketch}

\begin{remark}[the verification path]
\label{v4-arith:bzn:rem:verification-sketch}
A full verification paper (scope: 20--30 additional pages) would
proceed as follows.
\begin{enumerate}
\item \emph{Teleman--Antieau--Nikolaus computation of
$\mathrm{HH}_{\bullet}$ of the global Iwahori Hecke category via
THH$_{\bullet}$} of the Emerton--Gee stack regarded as a
prismatic-Nygaard-polarised derived scheme. This path goes through
the trace-methods framework of Nikolaus--Scholze, producing an
$E_{1}$-algebra in genuine $S^{1}$-equivariant spectra, whose
underlying object matches the factorization-homology model of
\Cref{v4-arith:bzn:thm:arith-fact-homology}.
\item \emph{MKR-dictionary match of the arithmetic loop space.}
From the topological analogy, the arithmetic free loop space
The full-target arithmetic loop object for
$\mathcal{X}^{\mathrm{Sel}}_{\bar\rho,\det,\mathcal{L}}$ corresponds
to the loop space of the $3$-manifold analogue of the Emerton--Gee
stack; Morishita's dictionary (arXiv:0904.3399 Ch.~3) reads
off the correct moduli of local-system pairs.
\item \emph{Comparison of Bhatt--Morrow--Scholze prismatic-Nygaard
trace with the $E_{1}$-algebra $\mathcal{O}(\mathrm{LocSys}^{\mathrm{Ar}})^{\mathrm{flat}}$.}
Both sides give an $E_{1}$-algebra in $S^{1}$-equivariant spectra;
the match is a prismatic-Nygaard compatibility, encoded as the
Bhatt--Lurie filtration of arXiv:2201.06120.
\end{enumerate}
The scope of this verification is deliberately smaller than the
40--60 page main proof because the verification is a
\emph{consistency check}, not a re-proof; it verifies that the
arithmetic Ben-Zvi--Nadler equivalence of
\Cref{v4-arith:bzn:thm:main-body} is consistent with independent
computations, not that an independent proof exists.
\end{remark}

\section{Residual Scope}
\label{v4-arith:bzn:residual}

The attack-and-repair run leaves F3 with genuine residual obligations.
Besides F1 and F2, the proof needs the corrected arithmetic BZN
comparison against the full Weil/$(\varphi,\Gamma)$ target, the
profinite formal-loop continuity theorem, and the flat-locus
pullback comparison. Thus F3.BZN$^{\mathrm{Ar}}$ remains a residual
of \Cref{v4-arith:cl:summary}.

\begin{remark}[what remains]
\label{v4-arith:bzn:rem:what-remains}
After \Cref{v4-arith:bzn:thm:main-body}, the open ingredients are:
\begin{itemize}
\item the local categorical comparison (F1.a) and
$\mathrm{L}_{\mathrm{ACD}}$ (F1's adelic categorical descent);
\item the mixed-characteristic local $E_{2}$ Hecke enhancement and
(F2.c) ($\infty$-categorical Tate restricted tensor product);
\item profinite formal-loop continuity for
$\mathrm{Map}(B\widehat{\mathbb{Z}},-)$;
\item the full Weil/$(\varphi,\Gamma)$ target comparison and
flat-locus pullback;
\item arithmetic Positselski (from P1's Step 3);
\item Koszul--Petersson compatibility (from P3).
\end{itemize}
The first four are number-field adaptations of function-field
geometric Langlands and arithmetic loop-space technology. The last
two are bar-cobar-structural obligations internal to the arithmetic
Heisenberg.
\end{remark}

\subsection{Scope: Archimedean Place}

\begin{remark}[archimedean contribution]
\label{v4-arith:bzn:rem:archimedean-contribution}
\Cref{v4-arith:bzn:thm:main-body} treats only the finite-place
inertia (\Cref{v4-arith:bzn:rem:archimedean}). The archimedean
contribution to $\mathcal{V}^{\mathrm{prim}}$ is the complex
Heisenberg Fock on the archimedean fibre (\Cref{v4-arith:thm:P3-resolution}
archimedean-place discussion), whose Hochschild trace is classical
Ben-Zvi--Nadler over $\mathbb{C}$. The arithmetic-to-classical
reduction at the archimedean place is exactly the classical
Ben-Zvi--Nadler theorem \Cref{v4-arith:bzn:thm:BZN-classical},
which is a theorem of 2009. No new work is needed for the
archimedean fibre.
\end{remark}

\subsection{Scope: Non-$GL_{2}$ Groups}

\begin{remark}[extension to $GL_{n}$]
\label{v4-arith:bzn:rem:GLn}
The machinery of this chapter extends to $G=GL_{n}$ for any $n\geq 1$,
since:
\begin{itemize}
\item The definition of $B\widehat{\mathbb{Z}}$ and arithmetic
free loop space is group-independent.
\item Local Emerton--Gee stacks are defined for all $n$ in
Emerton--Gee arXiv:1908.07185; the global spectral target must still
be Selmer-derived.
\item Fargues--Scholze local geometrization and spectral action are
available for $GL_{n}$ (arXiv:2102.13459 Thm.~I.9); the categorical
comparison remains a separate input.
\end{itemize}
The only $GL_{2}$-specific input is F1 ($\mathrm{L}_{\mathrm{ACD}}$
is currently formulated for $GL_{2}/\mathbb{Q}$), which is a
number-theoretic arithmetic-Langlands conjecture. With F1 upgraded
to $GL_{n}$, \Cref{v4-arith:bzn:thm:main-body} extends directly.
The rank-$2$ structure is not used essentially.
\end{remark}

\section{Summary}
\label{v4-arith:bzn:summary}

\begin{enumerate}
\item The arithmetic inertia circle $B\widehat{\mathbb{Z}}$ is
well-defined as a pro-\'etale derived stack, has pro-\'etale
fundamental group $\widehat{\mathbb{Z}}=\prod_{p}\mathbb{Z}_{p}$,
and is the arithmetic analogue of $S^{1}$ under the MKR dictionary
(\Cref{v4-arith:bzn:prop:MKR-inertia-circle}).
\item The arithmetic free loop space
$L^{\mathrm{Ar}}\mathcal{X}=\mathrm{Map}(B\widehat{\mathbb{Z}},\mathcal{X})$
is a derived pro-\'etale stack for any $n$-geometric derived
pro-\'etale stack $\mathcal{X}$
(\Cref{v4-arith:bzn:def:L-ar}, \Cref{v4-arith:bzn:rem:L-ar-existence}).
For $\mathcal{X}=BGL_{2}$, $L^{\mathrm{Ar}}BGL_{2}=\mathrm{LocSys}^{\mathrm{Ar}}_{GL_{2}}(B\widehat{\mathbb{Z}})$
(\Cref{v4-arith:bzn:prop:L-ar-BGL2}, \Cref{v4-arith:bzn:cor:L-ar-BGL2-is-LocSys}).
\item The Emerton--Gee stack has an expected comparison to the full
Weil/$(\varphi,\Gamma)$ flat target; its restriction to
$\mathrm{LocSys}^{\mathrm{Ar}}_{GL_{2}}(B\widehat{\mathbb{Z}})$ is a
forgetful procyclic shadow and is not known to be a closed immersion
(\Cref{v4-arith:bzn:prop:EG-to-LocSys},
\Cref{v4-arith:bzn:prop:EG-image-flat}).
\item Arithmetic factorization homology is conjectured to give
$\mathrm{HH}_{\bullet}(\mathcal{C})=\int_{B\widehat{\mathbb{Z}}}\mathcal{C}$
for any $E_{2}$-monoidal stable presentable $\infty$-category
$\mathcal{C}$ in pro-\'etale derived geometry
(\Cref{v4-arith:bzn:thm:arith-fact-homology}).
\item The core arithmetic HH identity
\Cref{v4-arith:bzn:thm:HH-IndCoh-arith} gives
$\mathrm{HH}_{\bullet}(\mathrm{IndCoh}_{\mathrm{Nilp}}(\mathcal{X}))=\mathcal{O}(L^{\mathrm{Ar}}\mathcal{X})^{\mathrm{flat}}$
for any Noetherian derived formal algebraic pro-\'etale stack
$\mathcal{X}$; this remains theorem-waiting for formal Noetherian
pro-\'etale stacks and requires the R1.formal.b/R1.formal.c inputs.
\item Dunn additivity gives the $E_{1}$-algebra structure on
the Hochschild trace (\Cref{v4-arith:bzn:prop:E1-on-HH}),
conditional on F2.
\item The main comparison \Cref{v4-arith:bzn:thm:main-body}
is conjectural and must target
$\mathrm{LocSys}^{W,\varphi\Gamma}_{GL_{2}}$, not the inertia-only
shadow.
\item F3 remains theorem-waiting after F1 and F2
(\Cref{v4-arith:bzn:thm:F3-is-theorem}); the residual is the
corrected BZN comparison, formal-loop compatibility, and full-target
pullback.
\item The independent verification path uses disjoint sources
(MKR topological dictionary, Antieau--Nikolaus arithmetic THH,
Teleman HKR, Bhatt--Morrow--Scholze prismatic Nygaard) and
provides a consistency check for the main theorem
(\Cref{v4-arith:bzn:prop:disjoint},
\Cref{v4-arith:bzn:rem:verification-sketch}).
\item The machinery extends to $GL_{n}$ for arbitrary $n\geq 1$
with only F1 as the group-specific input
(\Cref{v4-arith:bzn:rem:GLn}).
\end{enumerate}

The six-window presentation of $\mathcal{V}^{\mathrm{prim}}$
(Chapter~\ref{v4-arith:synthesis}) remains theorem-waiting. F1 and
F2 do not remove F3. The residual obstruction is now sharper:
construct the corrected arithmetic BZN equivalence over the full
Weil/$(\varphi,\Gamma)$ target and prove the profinite formal-loop
and flat-locus comparisons.
